\chapter{One LC Stage}
\label{ch:single-stage}

A single series $L$ feeding a shunt $C$ is the building block for everything
else. Get the transfer function, the resonant peak and the output impedance of
one stage right, and the ladders in later chapters are just more of the same algebra.

\section{Transfer function}

With an ideal source, one stage is $H = Z_C/(Z_C + Z_L)$ from
\zcref{eq:current-divider}. Put both parasitics in:
$Z_L = R_L + sL$ and $Z_C = \ESR + 1/(sC)$.

\begin{derivation}
	Multiply top and bottom by $sC$:
	\[
		H(s) = \frac{\ESR + \frac{1}{sC}}{\ESR + \frac{1}{sC} + R_L + sL}
		= \frac{1 + sC\,\ESR}{1 + sC\,\Rloop + s^2 LC},
		\qquad \Rloop = R_L + \ESR.
	\]
	Now match it term by term against the standard second-order template:
	\[
		H(s) = \frac{1 + s/\omega_z}{1 + \dfrac{s}{Q\omega_0} + \left(\dfrac{s}{\omega_0}\right)^2}.
	\]
	The $s^2$ term gives $\omega_0^2 = 1/(LC)$. The $s$ term gives
	\[
		\frac{1}{Q\omega_0} = C\Rloop
		\quad\Rightarrow\quad
		Q = \frac{1}{\omega_0 C \Rloop} = \frac{R_0}{\Rloop},
	\]
	using $1/(\omega_0 C) = \sqrt{L/C} = R_0$. The numerator gives the ESR zero
	$\omega_z = 1/(C\,\ESR)$.
\end{derivation}

With $x = f/f_0$, the magnitude is
\begin{equation}\label{eq:H-single-mag}
	\abs{H} = \frac{\sqrt{1 + (f/f_z)^2}}{\sqrt{(1 - x^2)^2 + (x/Q)^2}}.
\end{equation}

\section{The resonant peak}

At $x = 1$ the $(1 - x^2)$ term vanishes and only the damping term is left:
\begin{equation}\label{eq:H-peak}
	\abs{H(f_0)} = Q\sqrt{1 + (f_0/f_z)^2} \approx Q.
\end{equation}

\begin{idea}
	A filter \emph{amplifies} noise at its resonance, by a factor of $Q$.
	A stage with $Q = 25$ turns a noise line sitting at $f_0$ into one 28~dB
	louder. $Q$ is ``reactance over loss'': how many times the resonance amplifies.
\end{idea}

The exact peak of \zcref{eq:H-single-mag} (ignoring the ESR zero) sits slightly
below $f_0$, at $x = \sqrt{1 - 1/(2Q^2)}$, with height $Q/\sqrt{1 - 1/(4Q^2)}$.
For any $Q$ worth worrying about, that is $f_0$ and $Q$.

\section{Above resonance}
\label{sec:above-resonance}

\begin{derivation}
	Above resonance $\abs{Z_C} \ll \abs{Z_L}$. In the lossless case,
	\[
		\abs{H} = \frac{1}{x^2 - 1} \approx \frac{1}{x^2}
		= \frac{1}{\omega^2 LC} = \frac{\omega_0^2}{\omega^2}
		= \left(\frac{f_0}{f}\right)^2.
	\]
	In dB, the attenuation $A = 1/\abs{H}$ is
	\[
		A_{dB} = 40\log_{10}\frac{f}{f_0}.
	\]
	The approximation is off by the factor $x^2/(x^2 - 1)$, which is under
	0.1~dB once $f > 5f_0$.
\end{derivation}

\section{Output impedance peak}
\label{sec:zo-peak}

Looking back from the converter with the source shorted, $L$ and $C$ are in parallel:
\[
	\Zo = Z_L \parallel Z_C = \frac{Z_L Z_C}{Z_L + Z_C}.
\]

\begin{derivation}
	At $\omega_0$ the reactances cancel in the denominator, leaving only the losses:
	\[
		Z_L + Z_C = \Rloop + j\left(\omega_0 L - \frac{1}{\omega_0 C}\right) = \Rloop.
	\]
	The numerator is barely affected by the small losses:
	\[
		Z_L Z_C \approx (j\omega_0 L)\left(\frac{1}{j\omega_0 C}\right) = \frac{L}{C} = R_0^2.
	\]
	So
	\begin{equation}\label{eq:zo-peak}
		\abs{\Zo}_{peak} \approx \frac{R_0^2}{\Rloop} = QR_0.
	\end{equation}
	Physically: at $f_0$ the loop current is limited only by $\Rloop = \DCR + \ESR$.
\end{derivation}

\begin{idea}
	Low-loss parts make a \emph{worse} filter for stability. Halving the DCR
	doubles the $\Zo$ peak. That is why damping is added on purpose (\zcref{ch:damping}).
\end{idea}

\section{Parasitics: where each one matters}

\begin{itemize}
	\item Series parasitics set $Q$. Shunt parasitics make zeros.
	\item Each parasitic matters only near its own frequency.
	\item Series $R$ matters within about $\pm f_0/(2Q)$ of $f_0$.
	\item ESR matters approaching and above $f_z = 1/(2\pi C\,\ESR)$, where it
	      bends the $-40$~dB/decade slope toward $-20$.
\end{itemize}

\begin{examplebox}[label=ex:first-try]{a first-try LC}
	A natural first attempt for the example device is a single stage with
	$L = \qty{68}{\micro\henry}$ (\qty{0.12}{\ohm} DCR) and
	$C = \qty{6.8}{\micro\farad}$ (\qty{10}{\milli\ohm} ESR), on an ideal source.

	\textbf{Resonance.}
	\[
		f_0 = \frac{1}{2\pi\sqrt{(\qty{68}{\micro\henry})(\qty{6.8}{\micro\farad})}} = \qty{7.40}{\kHz},
		\qquad R_0 = \sqrt{\frac{68}{6.8}}\,\Omega = \qty{3.16}{\ohm}.
	\]

	\textbf{Output impedance.} $\Rloop = 0.12 + 0.01 = \qty{0.13}{\ohm}$, so
	\[
		Q = \frac{3.16}{0.13} = 24.3, \qquad
		\abs{\Zo}_{peak} \approx \frac{R_0^2}{\Rloop} = \frac{10}{0.13} = \qty{76.9}{\ohm}.
	\]
	An exact sweep gives \qty{77.0}{\ohm}, so the shortcut is excellent.

	\textbf{Stability.} Against $\abs{\Zin} = \qty{85.0}{\ohm}$ from \zcref{ex:zin}:
	\[
		\text{margin} = 20\log_{10}\frac{85.0}{77.0} = 0.9~\text{dB}.
	\]
	That fails the 10~dB target badly. This innocent-looking filter is within
	10\% of oscillating.

	\textbf{Attenuation.} At \qty{400}{\kHz} the asymptote gives
	$\abs{H} \approx (7.40/400)^2$, or $-69.3$~dB. The exact value is $-69.2$~dB.

	\textbf{Resonant peak.} $\abs{H}$ peaks at $20\log_{10} 24.3 = +27.7$~dB at
	\qty{7.40}{\kHz}. Any noise near \qty{7.4}{\kHz} comes out 24 times larger.
\end{examplebox}

% -------------------------------------------------------------------

\section{Sizing from the emission requirement}
\label{sec:requirement}

The attenuation requirement sets how low the filter's resonance must sit.
Use the worst case: lowest input voltage, full load, lowest switching frequency.
\[
	P_{in} = \frac{P_{out}}{\eta}, \qquad I_{avg} = \frac{P_{in}}{V_{min}}.
\]

\subsection{Bounding the noise current}

Without a measured spectrum, you still know the converter's input current is a
pulse train that never goes negative. That alone bounds every harmonic.

\begin{derivation}
	The complex Fourier coefficient of harmonic $k$ is
	\[
		c_k = \frac{1}{T}\int_0^T i(t)\,e^{-jk\omega t}\,dt.
	\]
	The amplitude of the real sinusoid at $k\omega$ is $\abs{I_k} = 2\abs{c_k}$.
	Since $i(t) \geq 0$ and $\abs{e^{-jk\omega t}} = 1$,
	\[
		\abs{I_k} = \frac{2}{T}\abs*{\int_0^T i(t)\,e^{-jk\omega t}\,dt}
		\leq \frac{2}{T}\int_0^T i(t)\,dt = 2I_{avg}.
	\]
\end{derivation}

\begin{derivation}
	The bound is tight for narrow pulses. For a rectangular pulse train of
	height $I_p$ and width $\tau$, the fundamental is
	\[
		\abs{I_1} = \frac{2I_p}{\pi}\sin\frac{\pi\tau}{T}.
	\]
	For $\tau \ll T$, $\sin(\pi\tau/T) \approx \pi\tau/T$, so
	\[
		\abs{I_1} \approx \frac{2I_p}{\pi}\cdot\frac{\pi\tau}{T} = 2I_p\frac{\tau}{T} = 2I_{avg}.
	\]
	So $2I_{avg}$ is the right conservative estimate.
\end{derivation}

Convert to the units of the limit:
\[
	I_{rms} = \frac{I_{pk}}{\sqrt{2}}, \qquad
	N_{\text{dB}\mu\text{A}} = 20\log_{10}\frac{I_{rms}}{\qty{1}{\micro\ampere}}.
\]

\subsection{Required attenuation}

\begin{equation}\label{eq:A-req}
	A_{req} = N_{\text{dB}\mu\text{A}}(f_{sw}) - \text{Limit}_{\text{dB}\mu\text{A}}(f_{sw}) + \text{margin}.
\end{equation}

Use the lowest switching frequency the converter reaches at your load and
line. A resonant (ZCS) converter's frequency moves with load, and attenuation
is worst at the low end.

\subsection{One stage or two}

Invert the asymptotes. One stage needs $40\log_{10}(f_{sw}/f_0) = A_{req}$:
\begin{equation}\label{eq:one-stage-size}
	f_0 = \frac{f_{sw}}{10^{A_{req}/40}}, \qquad L = \frac{1}{(2\pi f_0)^2 C}.
\end{equation}

Two stages need $40\log_{10}(f_{sw}^2/(f_{01}f_{02})) = A_{req}$, which fixes
only the \emph{product} of the stage resonances:
\begin{equation}\label{eq:two-stage-size}
	f_{01}f_{02} = \frac{f_{sw}^2}{10^{A_{req}/40}}.
\end{equation}

Its geometric mean $\sqrt{f_{01}f_{02}}$ is a useful ``typical'' stage
frequency, and it sits much higher than the single-stage $f_0$. That is the
whole case for two stages: higher resonances mean smaller parts.

\begin{idea}
	Budget penalties up front. Fold every known loss of attenuation into one number,
	valid where the asymptotes hold:
	\[
		A_{req}' = A_{req} + \sum \text{penalties}_{dB}
		\quad \text{(damping, cap bias, tolerance)}.
	\]
\end{idea}

Two handy dB facts: a factor of 2 is 6.02~dB and a factor of 10 is 20~dB.
Sanity check any $H$: as $f \to \infty$, $\abs{H}$ must go to 0.

\begin{examplebox}[label=ex:requirement]{the attenuation requirement}
	\textbf{Noise current.} At \qty{100}{\volt} and full load,
	$I_{avg} = \qty{117.6}{\watt}/\qty{100}{\volt} = \qty{1.176}{\ampere}$.
	The fundamental is bounded by
	\[
		\abs{I_1} \leq 2I_{avg} = \qty{2.35}{\ampere}~\text{pk}
		= \qty{1.66}{\ampere}~\text{rms}
		= 20\log_{10}(1.66\times 10^{6}) = 124.4~\text{\dBuA}.
	\]

	\textbf{Requirement.} From \zcref{eq:A-req}:
	\[
		A_{req} = 124.4 - 50 + 6 = 80.4~\text{dB}.
	\]
	The first-try filter of \zcref{ex:first-try} gives 69.2~dB, so it is
	11.2~dB short as well as unstable.

	\textbf{One stage.} From \zcref{eq:one-stage-size}:
	\[
		f_0 = \frac{\qty{400}{\kHz}}{10^{80.4/40}} = \qty{3.90}{\kHz},
		\qquad
		L = \frac{1}{(2\pi \cdot \qty{3.90}{\kHz})^2(\qty{6.8}{\micro\farad})} = \qty{244}{\micro\henry}.
	\]
	That is before any damping, which will cost more (\zcref{ex:single-damped}).

	\textbf{Two stages.} From \zcref{eq:two-stage-size}:
	\[
		f_{01}f_{02} = \frac{(\qty{400}{\kHz})^2}{10^{80.4/40}} = \qty{1.56e9}{\hertz\squared},
		\qquad \sqrt{f_{01}f_{02}} = \qty{39.5}{\kHz}.
	\]
	Each stage can resonate about ten times higher than the single stage, so
	the parts are much smaller.
\end{examplebox}
